\section{RELATED WORK}
Research relevant to this work spans CPS resilient recovery, reinforcement learning for sequential recovery, and trustworthy federated learning.

\subsection{Resilient Recovery of CPS}
CPS resilience concerns the ability to withstand disruption and restore acceptable physical operation after faults or attacks \cite{segovia2024survey,o2024towards,lu2024recovery}. Existing recovery studies use robust control, observer reconstruction, optimization, game theory, and learning-based control \cite{yuan2024incomplete,kim2022survey,niu2022verifying,yu2024intelligent,xin2025learning,zhong2024deep}. These methods establish important physical recovery principles, but many assume that the recovery controller or the information consumed by that controller remains trustworthy. Our work focuses on the harder case in which the physical plant is under attack while the distributed learning process used to adapt recovery policies can itself contain Byzantine participants.

A central distinction in the revised formulation is between \emph{learning quality} and \emph{physical safety}. Robust aggregation can reduce poisoned-model influence, but no statistical aggregation rule can guarantee that every learned action is safe under arbitrary poisoning. We therefore place a model-independent safety-recovery contract between the learned policy and the physical plant. This design makes safety conditional on explicit state/attack uncertainty coverage and shield feasibility rather than on successful convergence of the learning algorithm.

\subsection{Reinforcement Learning for CPS Recovery}
RL has been applied to CPS because it can optimize sequential decisions when accurate analytical models are incomplete \cite{olowononi2020resilient,nguyen2021deep}. Multi-agent RL extends this capability to distributed and coupled systems \cite{zhang2021physical,jin2025comprehensive}. Existing work has considered anomaly response, routing, distributed control, and deadline-aware decisions \cite{10105656,10393887,yang2024adt,liu2024deadline,wang2023multi}. These methods motivate the local actor--critic component used here.

However, a critic-estimated advantage is only a learning-utility signal. Under corrupted observations or poisoned critics, a high estimated advantage does not prove that a transition is physically safe. The revised AEFC therefore uses advantage only together with local recovery certificates for honest-client knowledge selection, and it never treats client-reported advantage as proof of trustworthiness at the server.

\subsection{Federated and Heterogeneous Learning}
Federated learning keeps raw data local but remains vulnerable to statistical heterogeneity, client drift, system heterogeneity, communication bottlenecks, privacy leakage, and malicious updates \cite{kairouz2021advances,li2020federated}. Personalized and representation-aligned methods such as FedARA, Fed-ADE, FedDAP, and F$^2$DC separate shared and local components or use alignment signals under client-specific distributions \cite{zhao2026fedara,park2026fedade,le2026feddap,wang2026f2dc}. These studies support the use of a shared recovery backbone with subsystem-specific observation and action adapters.

Our setting is heterogeneous FRL rather than classical vertical federated learning. Vertical FL assumes feature-partitioned records with aligned sample identities \cite{liu2024vertical}; here, clients interact with different local environments and generate different trajectories. Only the representation-compatible recovery backbone is federated, while local adapters absorb incompatible sensing and actuation dimensions.

Communication-efficient FL frequently relies on sparsification or compressed updates \cite{jia2025comprehensive,11090013}. Compression is not itself a security mechanism, and data-dependent sparse coordinates may still leak information. In this paper, Top-$k$ is used only to reduce communication. Its deterministic residual is analyzed explicitly instead of being assigned an unspecified contraction coefficient.

\subsection{Trustworthy and Byzantine-Resilient Federated Learning}
Secure FL studies address poisoning, backdoors, malicious clients, privacy leakage, and trusted execution through robust aggregation, filtering, client screening, TEEs, oblivious learning, and differential privacy \cite{miao2024esfl,miao2024asyfl,mu2024feddmc,zhao2026sxgb,kato2023olive}. Privacy-focused approaches such as DP-FedAdamW, FedMOP, and privacy-guaranteed personalized FL further demonstrate that data locality and sparsification are not substitutes for a quantifiable privacy mechanism \cite{liu2026dpfedadamw,zhao2026fedmop,patel2026pfdm}. We therefore make no formal privacy claim for the base AEFC mechanism and treat differential privacy as an optional extension.

Federated reinforcement learning has been studied for distributed control, malware mitigation, machine unlearning, predictive maintenance, and energy scheduling \cite{qi2021federated,10908884,10483278,10536634,qi2025federated}. CyberForce is especially relevant because it explicitly considers poisoning resilience \cite{10908884}. A2FC and related FRL architectures also show that advantage-based actors and federated critics are not new primitives in isolation. Accordingly, our novelty claim is not based on introducing advantage, Top-$k$, prototypes, or a federated critic individually.

The revised contribution is the security composition required by recovery integrity. Honest clients share only certificate-supported recovery knowledge; the server limits arbitrary Byzantine updates without trusting self-reported credibility; representation alignment is used to reduce honest heterogeneity rather than authenticate clients; and a robust safety-recovery shield prevents an arbitrarily poisoned learned proposal from directly violating the modeled physical safety set. The resulting guarantees are separated by layer: aggregation protects learning utility, while the execution certificate protects the physical plant.
